\section{Summary of Achievements}
\label{sec:conclusion_summary}

Chapter~2 reviewed the theoretical foundations underpinning the system: sequential recommendation via self-attention~\cite{kang2018self,sun2019bert4rec}, co-purchase graph embedding via Cleora~\cite{rychalska2021cleora}, approximate nearest-neighbour search via HNSW~\cite{malkov2018hnsw}, large-scale similarity indexing via FAISS~\cite{johnson2021faiss}, multilingual encoding via BGE-M3~\cite{chen2024bge}, and Reciprocal Rank Fusion for retrieval combination~\cite{cormack2009rrf}.
These methods were selected on the basis of empirical performance, latency characteristics, and compatibility with the existing data platform.

Chapter~3 analysed the limitations of the prior system.
The legacy GRU-SeqDQN recommendation module achieved HR@10 = 0.1031 — statistically indistinguishable from the random baseline of 0.1000 — due to the fixed-size GRU bottleneck and the intractable 3-million-item DQN action space.
The BLaIR text encoder produced Precision@10 $\leq 0.05$ on Vietnamese queries, making the system functionally unusable for the target multilingual user base.
The FAISS flat index incurred a P50 search latency of 612.19\,ms and a 45-second startup loading penalty, precluding interactive use.

Chapter~4 introduced four infrastructure changes to address these shortcomings.
BLaIR was replaced by BGE-M3, and the GRU-SeqDQN module was replaced by a dual-mode architecture comprising Pipeline~A (Cleora + BGE-M3) and Pipeline~B (DIF-SASRec).
A three-stage NLLB-200 translation pipeline extended query support to 19 non-English languages.
Parallel text and CLIP encoding was achieved via an \texttt{anyio} task group.
The flat FAISS index was replaced by a dual HNSW/flat strategy loaded via memory mapping, and a three-stage incremental Cleora refit pipeline was introduced to incorporate new interactions without restarting the server.

Chapter~5 evaluated these changes quantitatively.
Pipeline~A (Cleora + BGE-M3) achieves HR@10 = 0.9047 and Pipeline~B (DIF-SASRec) achieves HR@10 = 0.7745 on 100{,}000 held-out test users, both substantially above the Content Baseline (0.4346) and the prior GRU-SeqDQN approach (0.1031).
The DIF-SASRec result is stable across five random evaluation seeds (standard deviation $\sigma = 0.0024$).
Multilingual search performance reached Precision@10 = 0.55 on long Vietnamese queries under BGE-M3, compared to 0.05 under BLaIR.
Warm-cache end-to-end search latency was reduced from 1{,}102\,ms to 59\,ms.

Chapter~6 presented the deployed book discovery application, demonstrating the end-to-end integration of all components: the login gate, the multimodal search interface (text, image, and fused), the dual-panel recommendation view with the DIF-SASRec input-sequence strip, the real-time profile panel with the training loss sparkline, and the streaming Qwen2.5 book assistant.

---

\section{Evaluation of Objectives}
\label{sec:conclusion_objectives}

Table~\ref{tab:objectives} evaluates each objective stated in Section~\ref{sec:goals} against the experimental evidence.

\begin{table}[H]
\centering
\caption{Evaluation of project objectives against experimental evidence.}
\label{tab:objectives}
\begin{tabular}{p{5.8cm} p{2.2cm} p{5.5cm}}
\hline
\textbf{Objective} & \textbf{Status} & \textbf{Evidence} \\
\hline
Design a dual-mode recommendation pipeline (Pipeline~A + Pipeline~B) &
Met &
Table~\ref{tab:main_results}: Pipeline~A HR@10 = 0.9047, Pipeline~B HR@10 = 0.7745 \\

Achieve substantial improvement over the prior GRU-SeqDQN baseline &
Met &
HR@10 increased from 0.1031 (GRU-SeqDQN) to 0.9047 (Pipeline~A) and 0.7745 (Pipeline~B), Table~\ref{tab:main_results} \\

Enable multilingual semantic search with $\geq 5\times$ Precision@10 improvement on Vietnamese queries &
Met &
Precision@10 on long Vietnamese queries: 0.0500 (BLaIR) $\to$ 0.5500 (BGE-M3), an 11$\times$ improvement (Table~\ref{tab:query_retrieval}) \\

Reduce warm-cache end-to-end search latency below 100\,ms &
Met &
Latency reduced from 1{,}102\,ms to 59\,ms (warm cache), Table~\ref{tab:latency_stages} \\

Deploy an interactive application integrating all platform components &
Met &
Chapter~6 presents the deployed book discovery application combining multimodal search, dual-pipeline recommendations, online DIF-SASRec fine-tuning, and the Qwen2.5 LLM assistant \\
\hline
\end{tabular}
\end{table}

All five objectives are met.
The clearest improvement is in recommendation quality: the prior GRU-SeqDQN approach scored HR@10 = 0.1031, near random, while Pipeline~A reaches 0.9047 and Pipeline~B reaches 0.7745 on the same test set.

---

\section{Research Contributions}
\label{sec:contributions}

This project makes the following specific technical contributions:

\begin{enumerate}
    \item \textbf{Incremental Cleora hot-swap pipeline.}
    A three-stage fault-isolated pipeline (delta export → graph augmentation → atomic in-memory swap) enables the Cleora co-purchase embeddings to be updated from live interaction logs without server restart.
    The atomic swap guarantees that concurrent in-flight recommendation requests observe a consistent index state throughout the update.

    \item \textbf{DIF-SASRec decoupled attention with learnable fusion scalar.}
    The implementation extends the original DIF-SASRec formulation with a sigmoid-gated fusion scalar $\alpha$ that is initialised to favour the category stream ($\alpha_0 = 0.7$) and is learned jointly during pretraining.
    This allows the model to adapt the relative importance of content and category attention to the characteristics of the Amazon Books domain.

    \item \textbf{Min-max normalised DIF-SASRec intent scores for UI transparency.}
    Raw DIF-SASRec dot-product scores are unbounded and vary by orders of magnitude across users.
    The backend applies per-request min-max normalisation so that the top-ranked candidate always receives score 1.0 and all remaining candidates are proportional to their dot-product gap.
    This makes the per-card score badges on the \textit{You Might Like} panel directly interpretable by users.

    \item \textbf{Dual-panel recommendation interface with per-pipeline attribution.}
    The frontend exposes two recommendation panels that directly mirror the two backend pipelines, with per-card layer tags (\textit{Cleora\,+\,BGE-M3}, \textit{Cleora\,+\,CLIP}, \textit{DIF-SASRec}) and the DIF-SASRec input-sequence strip.
    This design makes the recommendation provenance visible to users and provides a basis for future A/B testing of pipeline variants.

    \item \textbf{training loss sparkline as training transparency.}
    A live SVG sparkline in the right-panel profile view plots the full per-user training loss history with an annotated phase boundary between pretraining and online fine-tuning.
    This provides immediate, per-interaction visual confirmation that the DIF-SASRec model has updated — a transparency mechanism with no equivalent in existing open-source recommendation systems.

    \item \textbf{BGE-M3-based semantic re-ranking for LLM grounding.}
    Rather than passing raw Wikipedia or Google Books extracts to the Qwen2.5 assistant, the grounding pipeline uses the same BGE-M3 encoder deployed throughout the system to select the top-5 plot-relevant sentences from the retrieved text.
    This suppresses author biography sentences and marketing language without requiring a separate re-ranking model.
\end{enumerate}

---

\section{Generalisation to Other Domains}
\label{sec:generalisation}

The system was designed with a modular architecture in which domain-specific and domain-agnostic components are cleanly separated.
Table~\ref{tab:generalisation} assesses the portability of each component to a hypothetical new domain.

\begin{table}[htbp]
\centering
\caption{Component portability to a new recommendation domain.}
\label{tab:generalisation}
\begin{tabular}{p{4.0cm} p{2.0cm} p{6.5cm}}
\hline
\textbf{Component} & \textbf{Portability} & \textbf{Notes} \\
\hline
FastAPI + React SPA framework & Drop-in & No domain coupling; only API schemas change \\
HNSW + FAISS retrieval        & Drop-in & Accepts any fixed-dimensional embedding \\
RRF fusion                    & Drop-in & Parameter-free; works with any pair of ranked lists \\
NLLB translation pipeline     & Drop-in & Supports 202 languages regardless of domain \\
Tantivy BM25 index            & Drop-in & Domain-agnostic text index \\
Interaction logging (Redis + MongoDB) & Drop-in & Schema generalises to any item-action event \\
BGE-M3 encoder                & Reusable & Multilingual; re-embedding required for new item catalogue \\
CLIP image index               & Reusable & Re-embedding required if image modality differs \\
AgentPool + checkpoint system  & Reusable & Architecture is independent of domain; pretraining required \\
DIF-SASRec pretrained weights  & Domain-specific & Pretrained on Amazon Books; retraining needed for new domain \\
Cleora co-purchase graph       & Domain-specific & Requires domain interaction logs to build the co-purchase graph \\
item\_metadata Parquet store   & Domain-specific & Schema (title, genre, image URL) maps naturally to many product domains \\
\hline
\end{tabular}
\end{table}

The retrieval infrastructure (HNSW, RRF, NLLB, Tantivy, Redis, MongoDB) requires no changes to support a new domain.
The encoder and model layers (BGE-M3, CLIP, DIF-SASRec) require re-embedding and retraining on domain-specific data, but the training code and serving infrastructure are reusable.
Only the Cleora graph is tightly coupled to the co-purchase interaction pattern; an analogous collaborative signal (e.g., co-view, co-play, co-purchase in another product domain) would enable direct reuse.

---

\section{Limitations}
\label{sec:conclusion_limitations}

Several limitations of the proposed system are identified.

\begin{itemize}
    \item \textbf{Incomplete catalogue coverage in Pipeline~A.}
    The Cleora co-purchase graph covers 375{,}280 items — 12.5\% of the full 3-million-item catalogue.
    Items outside the graph cannot be retrieved by Pipeline~A, contributing to the 2.6\% of test users (2{,}641) that neither pipeline serves (Section~\ref{sec:qualitative}).

    \item \textbf{Content veto over-filtering.}
    The cosine similarity threshold $\tau = 0.3$ in Pipeline~B occasionally filters out semantically legitimate candidates whose embeddings fall just below the threshold.
    This is particularly harmful for users with sparse histories, where the DIF-SASRec prediction may be the only available recommendation path.

    \item \textbf{HNSW recall degradation.}
    The HNSW index achieves a recall of 0.521 relative to exact flat search (Table~\ref{tab:hnsw_comparison}), meaning that 47.9\% of the nearest neighbours under the full index are not returned by the approximate index.
    The BM25 re-ranking stage partially compensates, but this gap may reduce recommendation quality for queries where exact top-$k$ proximity is critical.

    \item \textbf{Online fine-tuning latency on cold sessions.}
    The DIF-SASRec gradient step executes synchronously in the \texttt{POST /interact} request handler.
    For users with long interaction histories the checkpoint serialisation and deserialisation from disk introduces latency that grows proportionally with checkpoint size.

    \item \textbf{In-session cart state.}
    The shopping cart is held in React component state only and is not persisted to the backend.
    Cart contents are lost on page reload, preventing recovery of abandoned carts and limiting the cart signal from influencing Cleora graph updates across sessions.
\end{itemize}

---

\section{Future Work}
\label{sec:conclusion_future}

Several concrete directions for future improvement follow from the identified limitations.

\begin{itemize}
    \item \textbf{Expand the Cleora index to the full catalogue.}
    Extending the co-purchase graph from 375{,}280 to the full 3 million items would eliminate the catalogue boundary that causes Pipeline~A to fail for long-tail items.
    The incremental refit pipeline introduced in Section~\ref{sec:cleora_incremental} already supports partial augmentation; a batch re-embedding of the full catalogue with accumulated interaction data is the next step.

    \item \textbf{Adaptive content veto threshold.}
    Rather than a fixed threshold of $\tau = 0.3$, the veto could be relaxed to $\tau = 0.2$ or disabled entirely for users with fewer than 20 interactions, where over-filtering risk outweighs the benefit of semantic consistency enforcement.
    This directly addresses the sparse-user failure mode identified in Section~\ref{sec:qualitative}.

    \item \textbf{Asynchronous online fine-tuning.}
    Moving the DIF-SASRec gradient step from the synchronous \texttt{POST /interact} response path to a background task queue would decouple interaction logging latency from model update latency, allowing the training step to be batched and accelerated without impacting the user-facing response time.

    \item \textbf{Persistent cart and cross-session interaction signals.}
    Persisting cart events to MongoDB in the same manner as click and skip events would allow cart additions to contribute to Cleora graph updates across sessions, providing a higher-reward signal that currently disappears on page reload.

    \item \textbf{Extension to additional retrieval modalities.}
    The current retrieval pipeline combines text (BGE-M3) and visual (CLIP) signals.
    Incorporating structured metadata signals — genre taxonomy, publication year, reading-level classification — as additional FAISS or Tantivy filter axes would improve precision for queries that express constraints rather than pure semantic similarity (e.g., ``beginner-level machine learning books published after 2020'').
\end{itemize}
